\documentclass[12pt,a4paper]{report}

\usepackage[utf8]{inputenc}
\usepackage[T1]{fontenc}
\usepackage[brazil]{babel}
\usepackage{csquotes}
\usepackage{kpfonts}
\usepackage[a4paper,left=2.5cm,right=2.5cm,top=2.5cm,bottom=2.5cm]{geometry}
\usepackage{setspace}
\usepackage{microtype}
\usepackage{indentfirst}
\usepackage{fancyhdr}
\usepackage{amsmath,amssymb,amsthm,bm}

\usepackage{graphicx}
\usepackage{float}
\usepackage{caption}
\usepackage{subcaption}
\usepackage{booktabs}
\usepackage{multirow}
\usepackage{array}
\usepackage{pdflscape}
\usepackage{algorithm}
\usepackage{algpseudocode}
\usepackage{xcolor}
\usepackage{listings}
\usepackage[inline]{enumitem}
\usepackage{url}
\usepackage[style=numeric,backend=bibtex]{biblatex}
\addbibresource{refs.bib}
\usepackage[hidelinks]{hyperref}


\onehalfspacing
\setlength{\parindent}{1.25cm}
\setlength{\parskip}{0pt}
\setlength{\headheight}{14pt}
\captionsetup{font=small,labelfont=bf}
\graphicspath{{figures/}}

\renewcommand{\thesection}{\arabic{section}}
\renewcommand{\thesubsection}{\thesection.\arabic{subsection}}
\renewcommand{\thesubsubsection}{\thesubsection.\arabic{subsubsection}}
\renewcommand{\thefigure}{\arabic{figure}}
\renewcommand{\thetable}{\arabic{table}}
\renewcommand{\theequation}{\arabic{equation}}
\renewcommand{\thealgorithm}{\arabic{algorithm}}
\renewcommand{\contentsname}{Sumário}
\renewcommand{\bibname}{Referências}


\pagestyle{fancy}
\fancyhf{}
\fancyhead[L]{\small MC847 -- Aprendizado Não Supervisionado}
\fancyhead[R]{\small Trabalho Final}
\fancyfoot[C]{\thepage}
\renewcommand{\headrulewidth}{0.4pt}


\definecolor{codegreen}{rgb}{0,0.5,0}
\definecolor{codegray}{rgb}{0.5,0.5,0.5}
\definecolor{codepurple}{rgb}{0.58,0,0.82}
\definecolor{codeblue}{rgb}{0,0,0.65}
\definecolor{codebg}{rgb}{0.97,0.97,0.97}

\lstset{
  basicstyle=\ttfamily\footnotesize,
  backgroundcolor=\color{codebg},
  commentstyle=\color{codegreen},
  keywordstyle=\color{codeblue},
  stringstyle=\color{codepurple},
  numberstyle=\tiny\color{codegray},
  numbers=left,
  numbersep=5pt,
  frame=tb,
  breaklines=true,
  breakatwhitespace=false,
  showstringspaces=false,
  showtabs=false,
  tabsize=2,
  captionpos=b,
  keepspaces=true
}
\lstdefinestyle{python}{language=Python}
\lstdefinestyle{bash}{language=bash, numbers=none}


\newcommand{\todo}[1]{\textcolor{red}{\textbf{[TODO:} #1\textbf{]}}}


\newcommand{\includelogo}[1][0.25\textwidth]{%
  \IfFileExists{figures/Logo.png}{%
    \includegraphics[width=#1]{figures/Logo.png}%
  }{\IfFileExists{figures/Logo.pdf}{%
    \includegraphics[width=#1]{figures/Logo.pdf}%
  }{\IfFileExists{figures/Logo.jpg}{%
    \includegraphics[width=#1]{figures/Logo.jpg}%
  }{%
    \fbox{\parbox[c][2.5cm][c]{0.25\textwidth}{\centering\small
      Adicione \path{figures/Logo.png}}}%
  }}}%
}


\newcommand{\figplaceholder}[3]{%
  \begin{figure}[H]
    \centering
    \includegraphics[width=0.9\linewidth]{figures/#1}
    \caption{#2}
    \label{#3}
  \end{figure}
}

\begin{document}

\begin{titlepage}
  \centering
  \vspace*{0.5cm}

  \includelogo[0.22\textwidth]\\[1.2cm]

  {\large\bfseries MC847 -- Aprendizado Não Supervisionado}\\[0.3cm]
  {\large\bfseries Professor: Marcos Medeiros Raimundo}\\[1.4cm]

  {\large\bfseries Trabalho Final -- Relatório}\\[0.5cm]
  {\LARGE\bfseries Modelos de Recomendação e\\[0.2cm]
   Análise de Espaços Latentes}\\[2.5cm]

  {\large\bfseries Integrantes}\\[0.4cm]
  \begin{tabular*}{0.78\textwidth}{@{\extracolsep{\fill}} l r}
    Sara Beatriz da Silva Oliveira & 231288 \\[0.2cm]
    Integrante 2                    & RA 2   \\
  \end{tabular*}

  \vfill

  {\large Instituto de Computação}\\
  {\large Universidade Estadual de Campinas}\\[0.3cm]
  {\large Campinas, 2026}
\end{titlepage}

\clearpage
\pagenumbering{arabic}
\setcounter{page}{2}
\tableofcontents
\clearpage



\section{Problema, dados e protocolo}
% - O problema de recomendação como aprendizado transdutivo em matrizes esparsas, com a notação usada: 𝑅,𝐳𝑢,𝐰𝑖,Ω.
% - Análise dos dados: distribuição das notas, esparsidade, sinais de ausência não aleatória e perfil dos metadados.
% - Protocolo: divisão treino/teste, validação cruzada, intervalos de confiança e métricas (RMSE, MAE, NLL e a métrica de lista, com a definição de acerto e de candidatos).


\section{Modelos e suposições}
% - O que cada modelo supõe e como é ajustado:
%   - Baseline: vieses por mínimos quadrados regularizados.
%   - PCA com imputação: aproximação de posto 𝑘 e efeitos da imputação pela média.
%   - Fatoração de matrizes: perda nos dados observados, regularização 𝐿2 e SGD.
%   - PPCA: modelo gerador, algoritmo EM e posterior gaussiana.
%   - VAE: média e variância aprendidas por usuário, truque da reparametrização e ELBO.
% - Comparação entre modelos determinísticos e probabilísticos.


\section{Resultados}
% - Tabela com todos os modelos no teste, com intervalos de confiança: RMSE e MAE; NLL para PPCA e VAE; e a métrica de lista.
% - Comparação entre erro da nota (RMSE) e qualidade das listas.
% - Fatoração com e sem vieses.
% - PPCA e VAE linear com as mesmas suposições.
% - Efeito do decodificador não linear e sensibilidade a 𝛽, 𝜎2 ou 𝑘.
% - Convergência dos modelos, incluindo a curva em history do VAE.
% - Tempo de treino e de inferência.

\subsection{Etapa 1}

\subsection{Etapa 2}


\section{Análise dos Espaços Latentes}
% - Projeções 2D (PCA e t-SNE) dos quatro modelos.
% - Tabela com 𝑅2 (Ridge) e ROC-AUC (LogisticRegression) em dados não usados no ajuste, para os metadados de usuários e filmes.
% - Comparação de distâncias e vizinhanças entre os modelos.
% - Discussão: que informação cada espaço organiza melhor?


\section{Projeto}
% Pergunta, hipótese, experimento, previsão e o que refutaria a hipótese.
% O que foi implementado.
% Resultados: a hipótese foi confirmada ou refutada?
% O que o resultado ensina sobre o problema ou sobre o espaço latente.

\section{Conclusão}
% Principais achados.
% Limitações, como a abordagem transdutiva e o viés de seleção.
% Próximos passos.

\clearpage
\section{Referências}


\clearpage
\appendix
\renewcommand{\thesection}{\Alph{section}}
\setcounter{section}{0}
\section{Appendix}

\end{document}
